Geographic reliability is an important consideration when predictive models are evaluated using health-care payment data. This study examines prediction error across U.S. geographic areas using the 2023 Centers for Medicare \& Medicaid Services (CMS) Medicare durable medical equipment (DME) geography-and-service public-use file. After restricting the data to the 50 states and Washington, DC, the primary analysis included 36,528 aggregate records from 51 areas. We compared mean-based prediction baselines using Healthcare Common Procedure Coding System (HCPCS) codes alone and jointly with supplier rental classification. Leave-one-area-out evaluation was applied to actual and geographically standardized average Medicare payment targets. The HCPCS-only baseline yielded mean absolute errors (MAEs) of \$47.84 for actual payments and \$38.88 for standardized payments. Including rental classification reduced these MAEs to \$19.89 and \$10.30, respectively; the standardized target had lower area-level MAE in all 51 held-out areas. Among 916 HCPCS/rental groups present in at least ten areas, 805 exhibited lower coefficients of variation for standardized payments. Unseen HCPCS codes were infrequent but contributed disproportionately to absolute prediction error. These findings highlight the importance of equipment classification, target definition, and geographically disaggregated evaluation. They do not demonstrate that a model trained on standardized payments predicts actual Medicare expenditures more accurately.
